\begin{exercice*}[D'après DNB]
    Louise a téléchargé une liste de lecture sur son lecteur MP4.

    \begin{tabular}{|c|c|>{\centering\arraybackslash}m{0.25\linewidth}|}
        \hline
        Titre de la chanson & Nom de l'interprète & Durée de la chanson en secondes \\ \hline
        Mamatéou & Timaté & 232 \\ \hline
        La différence & Timaté & 211 \\ \hline
        Amazing & Timaté & 214 \\ \hline
        Tes racines & Timaté & 175 \\ \hline
        YoungBov & Hudad & 336 \\ \hline
        La ficelle & Maen & 191 \\ \hline
        Fou fou fou & Maen & 184 \\ \hline
        Nina & Maen & 217 \\ \hline
    \end{tabular}

    Louise décide d'utiliser la fonction « aléatoire » de son MP4. Cette fonction choisit au hasard une chanson parmi celles qui sont présentes dans la liste de lecture. Chaque chanson a la même probabilité d'être écoutée.

    \begin{enumerate}
        \item Détermine la probabilité que Louise écoute une chanson de Maen.
        \item Elle répète 25 fois l'utilisation de cette fonction et note à chaque fois le nom de l'interprète qu'elle a écouté. Les résultats qu'elle obtient sont notés dans le graphique ci-dessous.

        Quelle est la proportion de chansons de Maen écoutées ? Compare avec la question a.
    \end{enumerate}
\end{exercice*}
\begin{corrige}
    %\setcounter{partie}{0} % Pour s'assurer que le compteur de \partie est à zéro dans les corrigés
    %\phantom{rrr}    
    \begin{enumerate}
        \item La probabilité que Louise écoute une chanson de Maen est de $\dfrac{3}{8}$, car il y a 3 chansons de Maen sur les 8 chansons au total.
        \item La proportion de chansons de Maen écoutées est de $\dfrac{12}{25} = \num{0,48}$. Elle est donc supérieure à la probabilité calculée à la question a.
    \end{enumerate}
\end{corrige}

